\section{Translator, Assembler, Interpreter, dan Compiler}
\label{sec:ch03-translator}

\subsection{Keluarga Translator}

\begin{definisi}[Translator]
Translator adalah program yang menerjemahkan program sumber dari satu
bahasa ke bahasa lain. Keluaran translator dapat berupa bahasa
assembly, kode mesin, bytecode, atau bahasa tingkat tinggi lain.
\end{definisi}

Bahasa sumber dan bahasa target menentukan jenis translator. Assembler
menerjemahkan bahasa assembly ke kode mesin. Compiler menerjemahkan
bahasa tingkat tinggi ke bahasa tingkat rendah, biasanya assembly atau
kode mesin. Interpreter tidak menghasilkan kode target yang berdiri
sendiri; ia membaca program sumber dan langsung menjalankan aksinya.
Cross-compiler menghasilkan kode untuk platform yang berbeda dari
platform tempat compiler berjalan.

\begin{figure}[H]
\centering
\begin{tikzpicture}[node distance=5mm and 9mm]
  \node[phase, text width=3.0cm] (tr) {Translator};
  \node[phase, text width=2.8cm, below left=of tr] (as) {Assembler};
  \node[phase, text width=2.8cm, below=of tr] (co) {Compiler};
  \node[phase, text width=2.8cm, below right=of tr] (in) {Interpreter};
  \draw[arrow] (tr) -- (as);
  \draw[arrow] (tr) -- (co);
  \draw[arrow] (tr) -- (in);
\end{tikzpicture}
\caption{Keluarga translator dan tiga anggotanya yang utama}
\label{fig:ch03-translator}
\end{figure}

\subsection{Compiler versus Interpreter}

Compiler dan interpreter sama-sama membaca program sumber, tetapi
memiliki model eksekusi yang berbeda. Compiler menghasilkan program
target yang dapat dijalankan berulang kali tanpa kehadiran compiler.
Interpreter mengeksekusi program sumber secara langsung, biasanya
dengan bantuan representasi antara seperti bytecode.

\begin{table}[H]
\centering
\caption{Perbandingan compiler dan interpreter}
\label{tab:ch03-compiler-interpreter}
\begin{tabularx}{\textwidth}{@{}lXX@{}}
\toprule
\textbf{Aspek} & \textbf{Compiler} & \textbf{Interpreter} \\
\midrule
Keluaran & Kode target yang berdiri sendiri & Tidak ada kode target
permanen \\
Waktu eksekusi & Cepat setelah kompilasi & Lebih lambat karena
penerjemahan saat runtime \\
Waktu mulai & Lebih lama karena kompilasi & Lebih cepat \\
Portabilitas & Kode target terikat platform & Selama interpreter
tersedia \\
Analisis & Dapat melakukan analisis menyeluruh & Analisis terbatas
pada saat eksekusi \\
Contoh & GCC, Clang, javac & Python, Ruby, interpreter JavaScript \\
\bottomrule
\end{tabularx}
\end{table}

Batas antara compiler dan interpreter tidak selalu tegas. Java
mengompilasi sumber ke bytecode, lalu bytecode dijalankan oleh mesin
virtual yang dapat melakukan kompilasi just-in-time. Pendekatan
hybrid ini menggabungkan kelebihan keduanya.

\subsection{Cross-Compiler}

Cross-compiler berguna ketika program dikembangkan pada satu platform
tetapi dijalankan pada platform lain, misalnya kompilasi untuk
mikrokontroler dari komputer desktop. Compiler seperti ini harus
mengetahui konvensi pemanggilan, ukuran kata, dan tata letak memori
platform target. Pembahasan bootstrapping dan sejarah compiler tidak
menjadi fokus bab ini; keduanya cukup disebut sebagai konteks.

\begin{infobox}
Istilah \eng{compiler} berasal dari gagasan mengompilasi kumpulan
rutin menjadi satu program. Pada praktik modern, compiler juga
melakukan analisis, optimasi, dan pembangkitan kode, sehingga jauh
lebih kompleks daripada sekadar penggabung rutin.
\end{infobox}
